% La escalera de la certeza — decir exactamente lo que la evidencia aguanta.
%
% Diagrama propio en TikZ (sin licencias de terceros, vectorial, editable). Se
% tiñe con el color de la línea (milcGradeDark). Enseña el corazón ético de la
% fase de Evaluación: hay peldaños de afirmación, del más débil (opinión) al más
% fuerte (prueba), y la honestidad consiste en pararse en el que tu evidencia
% sostiene ---ni más abajo (falsa modestia) ni más arriba (exageración).
%
% Se incrusta vía \guiaDiagrama desde el bloque `recursos:` del YAML.
\newcommand{\peldano}[5]{% y · relleno · color de texto · afirmación · ejemplo
  \fill[#2] (0,#1) rectangle (6.6,{#1+0.82});
  \draw[milcNegro!45,line width=.4pt] (0,#1) rectangle (6.6,{#1+0.82});
  \node[anchor=west,font=\scriptsize,text=#3] at (0.18,{#1+0.55}) {\textbf{#4}};
  \node[anchor=west,font=\scriptsize\itshape,text=#3] at (0.18,{#1+0.24}) {#5};
}
\begin{tikzpicture}[scale=1,font=\sffamily]

  % Peldaños: de abajo (débil) a arriba (fuerte).
  \peldano{0}{milcGradeDark!12}{milcNegro}{creo / me pareció}{opinión, sin dato}
  \peldano{1}{milcGradeDark!32}{milcNegro}{observé / medí}{un dato tomado con método}
  \peldano{2}{milcGradeDark!55}{white}{mis datos sugieren}{un patrón, aún por confirmar}
  \peldano{3}{milcGradeDark!78}{white}{hay buena evidencia de}{varias medidas coinciden}
  \peldano{4}{milcGradeDark}{white}{está demostrado}{comprobado y repetido por otros}

  % Flecha de intensidad.
  \draw[->,milcNegro!70,line width=.9pt] (-0.45,0) -- (-0.45,4.82)
    node[midway,rotate=90,anchor=south,font=\scriptsize\bfseries,yshift=2pt] {más fuerte la afirmación};

  % Dónde puede pararse honestamente un proyecto del semillero.
  \draw[decorate,decoration={brace,amplitude=5pt},milcTurquesa,line width=.8pt]
    (6.85,1.0) -- (6.85,2.82);
  \node[anchor=west,font=\scriptsize\bfseries,milcTurquesa,align=left,text width=4.1cm]
    at (7.05,1.9)
    {aquí puede pararse tu proyecto (3 noches, 1 fotómetro): «observé» y «mis datos sugieren»};

  \node[anchor=west,font=\scriptsize,milcNegro!60,align=left,text width=4.0cm]
    at (7.1,4.35)
    {subir aquí sin más evidencia es \textbf{exagerar}};
  \draw[->,milcNegro!55] (8.4,4.08) -- (6.75,4.35);

  % Moraleja.
  \node[font=\scriptsize\itshape,milcNegro!75,align=center,text width=11cm]
    at (5.3,-0.95)
    {La honestidad no es decir \emph{menos} de lo que hallaste; es decir
     \textbf{exactamente} lo que tu evidencia aguanta. Cuidar el conocimiento
     empieza por no subir un peldaño que no te has ganado.};

\end{tikzpicture}
